\section{Method}
\label{sec:method}

\subsection{Problem formulation}

The prediction unit is one receptor structure, one ligand chemical entity, and
one explicitly specified binding pocket. The model receives receptor
coordinates and atom/residue identities, the ligand molecular graph and
stereochemistry, and a pocket center. It does not receive the crystallographic
ligand coordinates, pose RMSD, benchmark identity, or any feature computed
from a downstream evaluation. The output is a set of sampled all-atom ligand
poses and a separately declared pose-selection score. \method is a pose
generator and ranker; it is not trained as an affinity predictor or a blind
pocket detector.

Let the ligand contain \(N\) heavy atoms partitioned into \(K\) rigid
fragments. Atom \(i\) has a fragment assignment \(g(i)\) and a fixed
fragment-local coordinate \(\mathbf{y}_i\). A ligand pose is parameterized by
one translation and one proper rotation per fragment,
\[
    \mathbf{z}
    =
    \left\{(\mathbf{T}_k,\mathbf{R}_k)
    \mid \mathbf{T}_k\in\R^3,\,
    \mathbf{R}_k\in\SO(3)\right\}_{k=1}^{K},
\]
and the corresponding global atom coordinate is
\begin{equation}
    \mathbf{x}_i
    =
    \mathbf{R}_{g(i)}\mathbf{y}_i+\mathbf{T}_{g(i)}.
    \label{eq:atom_reconstruction}
\end{equation}
This representation preserves the internal geometry of each fragment while
allowing relative fragment motions to express ligand flexibility.

\subsection{Rigid-fragment decomposition}
\label{sec:fragmentation}

Ligands are sanitized and featurized with RDKit
\citep{landrum2024rdkit}. We identify cuttable bonds as heavy-atom single
bonds that are outside rings, are not planar conjugated bonds such as
amide-like linkages, and have endpoints with heavy-atom degree greater than
one. Candidate bonds are processed in deterministic order. A cut is accepted
only when each endpoint retains at least one non-cut heavy-atom bond, which
prevents isolated atoms along simple chains. Any remaining singleton component
is merged into its largest adjacent fragment. The resulting connected
components define \(g(i)\).

For each retained cut bond, we preserve the fragment adjacency and add local
cross-fragment triangulation edges around the cut-bond atoms. Their reference
distances provide a geometry signal without fixing the torsion angle. Fragment
centers are the geometric means of their member atoms, and local coordinates
are formed by subtracting the corresponding center. Single-atom fragments,
when unavoidable, have no supervised rotational degree of freedom.

\begin{figure}[t]
    \centering
    \includegraphics[width=\textwidth]{fragmentation_examples.png}
    \caption{Rigid-fragment representation. Cuttable rotatable bonds are
    highlighted in red, and atoms belonging to the same connected rigid
    component share a color. The model predicts one translation and one proper
    rotation per retained fragment while preserving atom-level chemistry and
    fragment adjacency.}
    \label{fig:fragmentation}
\end{figure}

\subsection{Pocket-conditioned heterogeneous graph}
\label{sec:graph}

The receptor is cropped in a residue-aware manner around the supplied pocket
center. During training, the current configuration samples the crop radius
uniformly from \(6\) to \(12\,\angstrom\) and applies Gaussian center jitter
with standard deviation \(2\,\angstrom\). The matched confidence and external
redocking protocol uses a fixed \(10\,\angstrom\) crop without center jitter.
At prospective inference, the pocket center must be supplied independently of
the unknown target pose.

We construct a single heterogeneous graph with four node types:
ligand atoms, ligand fragments, protein atoms, and protein-residue virtual
nodes. Ligand atom features encode element, formal charge, aromaticity,
hybridization, valence, chirality, ring membership, and pharmacophore flags.
Protein atom features encode residue--atom identity, backbone and metal
indicators, and protein-side pharmacophore flags. Residue virtual nodes encode
residue identity and provide a coarser structural communication path.

Ten typed edge relations represent ligand covalent bonds, local
triangulation, cut bonds, atom-to-fragment ownership, fragment adjacency,
protein covalent bonds, atom-to-residue membership, residue connectivity,
residue-to-fragment context, and dynamic protein-to-ligand contacts.
At flow time \(t\), the graph is
\(\mathcal{G}_t=(\mathcal{V},\mathcal{E}_{\mathrm{static}}\cup
\mathcal{E}_{\mathrm{contact},t})\). Protein--ligand atom contact edges are
rebuilt bidirectionally at every forward pass:
\begin{equation}
    \mathcal{E}_{\mathrm{contact},t}
    =
    \left\{
    (p,i),(i,p)
    \,\middle|\,
    p\in\mathcal{V}_{\mathrm{prot\text{-}atom}},
    i\in\mathcal{V}_{\mathrm{lig\text{-}atom}},
    \|\mathbf{x}_{p}-\mathbf{x}_{i,t}\|_2\leq 5\,\angstrom
    \right\}.
    \label{eq:dynamic_contacts}
\end{equation}
Edges are created only within the same complex in a minibatch.

For an edge \(j\rightarrow i\), define
\(\mathbf{r}_{ji}=\mathbf{x}_i-\mathbf{x}_j\) and
\(d_{ji}=\|\mathbf{r}_{ji}\|_2\). Edge attributes combine a 32-component
radial basis expansion with relation, bond, fragment-topology, local-frame,
pair-chemistry, endpoint-state, and time features. Their exact construction is
given in Appendix~\ref{app:edge_features}.

\subsection{Flow matching on fragment poses}
\label{sec:flow}

We center coordinates at the pocket and define a factorized prior over
fragment poses. For each fragment \(k\),
\begin{equation}
    \mathbf{T}_{k,0}
    \sim
    \mathcal{N}(\mathbf{0},\sigma^2\mathbf{I}),
    \qquad
    \mathbf{R}_{k,0}
    \sim
    \mathrm{Uniform}(\SO(3)).
    \label{eq:pose_prior}
\end{equation}
Single-atom fragment rotations are fixed to identity. During training,
\(\sigma\in\{0.5,1,2,3,4\}\,\angstrom\) is sampled with probabilities
\(\{0.10,0.25,0.30,0.25,0.10\}\), and \(\log\sigma\) is embedded along with
time. A global random proper rotation is applied to the ligand representation
as augmentation.

Given prior and target fragment poses
\((\mathbf{T}_{k,0},\mathbf{R}_{k,0})\) and
\((\mathbf{T}_{k,1},\mathbf{R}_{k,1})\), respectively, the default objective
uses a linear translation path and shortest-path spherical interpolation for
rotation:
\begin{align}
    \mathbf{T}_{k,t}
    &=
    (1-t)\mathbf{T}_{k,0}
    +t\mathbf{T}_{k,1},
    \label{eq:translation_path}\\
    \mathbf{R}_{k,t}
    &=
    \operatorname{Exp}(t\boldsymbol{\omega}^{\star}_k)
    \mathbf{R}_{k,0},
    \label{eq:rotation_path}\\
    \boldsymbol{\omega}^{\star}_k
    &=
    \operatorname{Log}
    \left(\mathbf{R}_{k,1}\mathbf{R}_{k,0}^{-1}\right).
\end{align}
The world-frame conditional velocity targets are therefore
\begin{equation}
    \mathbf{v}^{\star}_k
    =
    \mathbf{T}_{k,1}-\mathbf{T}_{k,0},
    \qquad
    \boldsymbol{\omega}^{\star}_k
    =
    \operatorname{Log}
    \left(\mathbf{R}_{k,1}\mathbf{R}_{k,0}^{-1}\right).
    \label{eq:flow_targets}
\end{equation}
Following the flow-matching formulation \citep{lipman2023flow}, a time-
conditioned network regresses these vector fields without integrating an ODE
during training. Time is sampled from
\[
  p(t)=0.02\,\mathcal{U}(0,1)
  +0.98\,\operatorname{sigmoid}
  \left(\mathcal{N}(0.8,1.7^2)\right),
\]
which emphasizes late refinement while retaining support near both endpoints.

\subsection{\texorpdfstring{\(\SE(3)\)}{SE(3)}-equivariant interaction network}
\label{sec:equivariant_network}

All node types share one six-layer equivariant trunk. Scalar chemical and
topological features are embedded into 384 channels. The full hidden state
uses the irreducible-representation layout
\[
    384\times 0e
    \;+\;
    32\times 1o
    \;+\;
    32\times 1e
    \;+\;
    16\times 2e
    \;+\;
    16\times 2o.
\]
Thus each node has 736 real-valued coordinates arranged into typed irrep
blocks rather than treated as an unstructured vector.

For complex \(b\), flow time and prior scale are embedded as
\begin{equation}
    \boldsymbol{\eta}_b
    =
    \operatorname{MLP}_{t}(\operatorname{PE}(t_b))
    +
    \operatorname{MLP}_{\sigma}
    \left(\operatorname{PE}(\log\sigma_b)\right)
    \in\R^{128}.
    \label{eq:time_sigma_embedding}
\end{equation}
Chemical and topological features initialize the \(0e\) block. The initial
\(1o\) channels are gated displacements from the graph center, with an
additional zero-initialized linear mixture of the current fragment rotation
columns. The \(1e\), \(2e\), and \(2o\) blocks start at zero.

For edge \(j\rightarrow i\), let \(\mathbf{e}_{ji}^{(\ell)}\) be the invariant
1020-dimensional descriptor and
\(\bar{\mathbf{h}}_j^{(\ell)}
=\operatorname{RMSNorm}(\mathbf{h}_j^{(\ell)})\). Each layer applies
dual radial channel scaling, a fully connected tensor product with normalized
spherical harmonics through \(l_{\max}=2\), and a direction-preserving
equivariant activation:
\begin{equation}
    \mathbf{m}_{ji}^{(\ell)}
    =
    \operatorname{EqAct}
    \left[
      \boldsymbol{\alpha}_{ji}^{\mathrm{out}}
      \odot
      \operatorname{TP}^{(\ell)}
      \left(
        \boldsymbol{\alpha}_{ji}^{\mathrm{in}}
        \odot\bar{\mathbf{h}}_j^{(\ell)},
        \mathbf{Y}_{0:2}(\widehat{\mathbf{r}}_{ji})
      \right)
    \right].
    \label{eq:tp_message}
\end{equation}
Scalar and non-scalar gates, including a learned edge-type-specific distance
decay, produce a normalized neighborhood aggregate
\(\mathbf{q}_i^{(\ell)}
=\operatorname{GateMean}_{j\in\mathcal{N}(i)}
(\mathbf{m}_{ji}^{(\ell)})\). The residual and time-conditioned update is
\begin{equation}
    \mathbf{h}_i^{(\ell+1)}
    =
    \operatorname{EqAdaLN}
    \left(
      \mathbf{h}_i^{(\ell)}
      +
      \operatorname{EqBlock}^{(\ell)}
      (\mathbf{q}_i^{(\ell)}),
      \boldsymbol{\eta}_{b(i)}
    \right).
    \label{eq:compact_layer_update}
\end{equation}
Appendix~\ref{app:architecture_details} gives the exact initialization, edge
descriptor, irrep normalization, radial gates, aggregation denominators, and
adaptive normalization.

Because coordinates enter only through equivariant geometric quantities and
all non-scalar operations respect their irreducible representations, scalar
features are invariant and predicted vectors are equivariant under a joint
proper rotation and translation of the complex. In particular, for
\(\mathbf{Q}\in\SO(3)\) and translation \(\mathbf{a}\),
\[
    \widehat{\mathbf{v}}(\mathbf{Q}\mathbf{x}+\mathbf{a})
    =
    \mathbf{Q}\widehat{\mathbf{v}}(\mathbf{x}),
    \qquad
    \widehat{\boldsymbol{\omega}}
    (\mathbf{Q}\mathbf{x}+\mathbf{a})
    =
    \mathbf{Q}\widehat{\boldsymbol{\omega}}(\mathbf{x}).
\]

\subsection{Atom-force head and Newton--Euler aggregation}
\label{sec:newton_euler}

The output head selects ligand-atom states and predicts one polar vector
\(\mathbf{f}_i\in\R^3\) per ligand atom. A linear equivariant path is augmented
with a self tensor product so that higher-order directional information can
reach the \(1o\) output:
\begin{align}
    \widetilde{\mathbf{h}}_i
    &=
    \operatorname{EqAct}
    (L_{\mathrm{pre}}\mathbf{h}_i^{(6)}),\\
    \mathbf{f}_i
    &=
    L_{1o}\widetilde{\mathbf{h}}_i
    +
    \operatorname{TP}_{1o}
    (\widetilde{\mathbf{h}}_i,\widetilde{\mathbf{h}}_i).
    \label{eq:atom_force_head}
\end{align}
The self-tensor-product branch is zero-initialized. The vectors are
``force-like'' equivariant readouts rather than gradients of a learned energy.
Rather than directly regressing two independent vectors for every fragment,
we convert them into rigid-body velocities. For fragment \(k\), with member set
\(\mathcal{A}_k\), center \(\mathbf{T}_k\), and lever arm
\(\mathbf{r}_i=\mathbf{x}_i-\mathbf{T}_k\),
\begin{align}
    \widehat{\mathbf{v}}_k
    &=
    \frac{1}{|\mathcal{A}_k|}
    \sum_{i\in\mathcal{A}_k}\mathbf{f}_i,
    \label{eq:newton_translation}\\
    \boldsymbol{\tau}_k
    &=
    \sum_{i\in\mathcal{A}_k}
    \mathbf{r}_i\times\mathbf{f}_i,\\
    \mathbf{I}_k
    &=
    \sum_{i\in\mathcal{A}_k}
    \left(
      \|\mathbf{r}_i\|_2^2\mathbf{I}_3
      -\mathbf{r}_i\mathbf{r}_i^\top
    \right),\\
    \widehat{\boldsymbol{\omega}}_k
    &=
    \mathbf{I}_k^{+}\boldsymbol{\tau}_k.
    \label{eq:newton_rotation}
\end{align}
We compute the pseudoinverse in the inertia eigenspace. Eigenmodes below 1\%
of the largest eigenvalue are treated as unobservable, so single-atom
fragments have no angular output and nearly singular modes of collinear
fragments are excluded. The corresponding projection \(\mathbf{P}_k\) is
applied to the angular target in the primary flow loss. The exact eigenspace
construction is given in Appendix~\ref{app:newton_euler_details}.

\subsection{Training objective}
\label{sec:objective}

Let
\(\mathcal{K}_{\mathrm{rot}}=\{k:|\mathcal{A}_k|>1\}\). The primary
translation and observable-rotation losses are
\begin{align}
    \mathcal{L}_{v}
    &=
    \frac{1}{3K}
    \sum_{k=1}^{K}
    \|\widehat{\mathbf{v}}_k-\mathbf{v}^{\star}_k\|_2^2,\\
    \mathcal{L}_{\omega}
    &=
    \frac{1}{3|\mathcal{K}_{\mathrm{rot}}|}
    \sum_{k\in\mathcal{K}_{\mathrm{rot}}}
    \|
      \widehat{\boldsymbol{\omega}}_k
      -\mathbf{P}_k\boldsymbol{\omega}^{\star}_k
    \|_2^2.
    \label{eq:primary_flow_losses}
\end{align}
Empty rotational sets contribute zero. The active configuration uses uniform
time weighting.

An atom-level auxiliary term compares the rigid-body atom velocities
\(\mathbf{u}_i=\mathbf{v}_{g(i)}
+\boldsymbol{\omega}_{g(i)}\times\mathbf{r}_i\).
A distance-geometry term extrapolates one predicted step to \(t=1\) and
penalizes inter-fragment atom-pair distance errors, weighted by \(t^2\).
Appendix~\ref{app:loss_details} gives both reductions explicitly. The total
objective is
\begin{equation}
    \mathcal{L}
    =
    \mathcal{L}_{v}
    +8\mathcal{L}_{\omega}
    +0.3\mathcal{L}_{\mathrm{atom}}
    +3.0\mathcal{L}_{\mathrm{DG}}.
    \label{eq:total_loss}
\end{equation}
All reductions use explicit valid-item counts so that distributed training on
unequal graph sizes matches a global masked mean.

\subsection{ODE sampling}
\label{sec:sampling}

At inference, independently sampled fragment priors are integrated from
\(t=0\) to \(t=1\). The default benchmark uses 25 Euler steps on the
late-biased grid
\begin{equation}
    t_s
    =
    1-\left(1-\frac{s}{S}\right)^3,
    \qquad s=0,\ldots,S,\quad S=25.
\end{equation}
At every step the ligand atom coordinates are reconstructed using
Eq.~\eqref{eq:atom_reconstruction}, dynamic contacts are rebuilt, and the
network predicts fragment velocities. Translations use an Euler update, while
rotations use left multiplication,
\[
    \mathbf{R}_{k,t+\Delta t}
    =
    \operatorname{Exp}
    \left(\Delta t\,
    \widehat{\boldsymbol{\omega}}_{k,t}\right)
    \mathbf{R}_{k,t}.
\]
The benchmark configuration is deterministic conditional on the complex ID,
global seed, and candidate index.

\subsection{Pose-confidence model}
\label{sec:confidence}

Sampling and top-1 selection are treated as separate decisions. The frozen
docking model first produces a set of \(M=80\) poses. For each candidate, a
second network initializes the same heterogeneous graph and adds the frozen
docking model's terminal ligand-atom and fragment irreps through a learned
gate initialized to 0.25. Four additional equivariant interaction layers then
process the full candidate complex.

The confidence network converts every terminal irrep state to an invariant
descriptor by concatenating its \(0e\) scalars with the norms of all
non-scalar channels. For each ligand atom, this descriptor is augmented with
explicit protein-contact features: minimum-distance RBFs, contact counts at
\(2\), \(3.5\), and \(5\,\angstrom\), a soft contact count, and
distance-weighted protein flags. An atom head predicts
\(\log(1+\)displacement\()\) and a \(2\,\angstrom\) success logit.

The deployed \texttt{global\_contact\_attention} readout combines local
contact attention with mean--maximum pooling over each of the four node types.
The final pose head is
\begin{equation}
    [\widehat{\ell}^{\,\mathrm{pose}}_m,
      \widehat{o}^{\,\mathrm{pose}}_m]
    =
    \operatorname{MLP}_{\mathrm{pose}}
    \left(
      \operatorname{LN}
      [\mathbf{r}^{\mathrm{global}}_m,
       \mathbf{r}^{\mathrm{contact}}_m]
    \right).
    \label{eq:confidence_pose_head}
\end{equation}
It predicts pose-level \(\log(1+\mathrm{RMSD})\) and a \(2\,\angstrom\)
success logit. Appendix~\ref{app:confidence_details} gives the invariant
features, attention, pooling, and output transformation explicitly.

For a pose set, the confidence objective combines atom and pose Huber
regression, binary success classification at \(2\,\angstrom\), pairwise
ranking, and a success-listwise term:
\begin{align}
    \mathcal{L}_{\mathrm{conf}}
    ={}&
    0.2\mathcal{L}_{\mathrm{atom\text{-}reg}}
    +0.2\mathcal{L}_{\mathrm{atom\text{-}BCE}}
    +0.3\mathcal{L}_{\mathrm{pose\text{-}reg}}
    +0.4\mathcal{L}_{\mathrm{pose\text{-}BCE}}
    \nonumber\\
    &+
    0.1\mathcal{L}_{\mathrm{rank}}
    +1.0\mathcal{L}_{\mathrm{success\text{-}listwise}}.
    \label{eq:confidence_loss}
\end{align}
The exact targets and pointwise, pairwise, and listwise reductions are given
in Appendix~\ref{app:confidence_loss_details}. The deployable default selects
the pose with minimum predicted RMSD. These predictions are used only for
within-set ranking and are not claimed to be calibrated RMSD estimates across
datasets or sampling distributions. A historical composite selector is
retained solely for exact compatibility analyses and is reported separately.
